%%%%%%%%%%%%%%%%%%%%%%%%%%%%%%%%%%%%%%%%%%%%%%%%%%%%%%%%%%%%%%%%%%%%%%%%%%%%%%%%%%
\begin{frame}[fragile]\frametitle{}
\begin{center}
{\Large Prompt Engineering for Educators}

\end{center}
\end{frame}


%%%%%%%%%%%%%%%%%%%%%%%%%%%%%%%%%%%%%%%%%%%%%%%%%%%%%%%%%%%
\begin{frame}[fragile]\frametitle{Why This Matters for You}

\begin{columns}
    \begin{column}[T]{0.55\linewidth}
    You do not need to be a programmer to benefit from Prompt Engineering.

    \medskip
    \textbf{These are real tasks you do every week:}
    \begin{itemize}
        \item Planning lessons or training sessions
        \item Giving feedback on student or colleague work
        \item Organising your own thoughts before a meeting
    \end{itemize}
    \end{column}
    \begin{column}[T]{0.45\linewidth}
    \begin{center}
    \renewcommand{\arraystretch}{1.4}
    \begin{tabular}{|l|c|}
    \hline
    \textbf{Task} & \textbf{Time Saved} \\
    \hline
    Lesson plan & 60\% \\
    Feedback writing & 75\% \\
    \hline
    \end{tabular}
    \end{center}
    \medskip
    {\tiny Estimates based on typical knowledge worker time studies.}
	
	    \textbf{A well-crafted prompt can cut each of these tasks from 30 minutes to 3 minutes.}
    \end{column}
\end{columns}

\end{frame}


%%%%%%%%%%%%%%%%%%%%%%%%%%%%%%%%%%%%%%%%%%%%%%%%%%%%%%%%%%%
\begin{frame}[fragile]\frametitle{Creating a Lesson Plan (For Teachers)}

    \begin{itemize}
        \item \textbf{Situation:} You need to plan a 45-minute lesson on a topic for a specific age group, today.
        \item \textbf{Result:} A complete, structured, ready-to-use lesson plan in under 30 seconds. Adapt the output to your classroom's specific needs.
    \end{itemize}
		
\begin{lstlisting}[language=HTML]
Prompt:
Create a detailed 45-minute lesson plan for:
- Topic: Introduction to Fractions
- Grade: 5th standard (age 10-11)
- Prior knowledge: students know division basics
- Learning objective: understand what a fraction means 
  and compare simple fractions (1/2, 1/4, 3/4)

Include:
- 5-minute warm-up activity (no materials needed)
- Main concept explanation with a real-life analogy
- A hands-on activity using items in a classroom
- 5 practice problems (easy to hard)
- A 2-minute closing question to check understanding
- One common misconception to watch out for
\end{lstlisting}



\end{frame}


%%%%%%%%%%%%%%%%%%%%%%%%%%%%%%%%%%%%%%%%%%%%%%%%%%%%%%%%%%%
\begin{frame}[fragile]\frametitle{Generating Quiz and Exam Questions}

    \begin{itemize}
        \item \textbf{Situation:} You need to create 20 questions for an assessment. Writing good questions is time-consuming, especially writing good \emph{wrong} answers for multiple-choice.
        \item \textbf{Why ``plausible wrong options'' matters:} Vague wrong answers make questions trivially easy. Asking the AI to use common misconceptions creates genuinely diagnostic questions.
    \end{itemize}
		
\begin{lstlisting}[language=HTML]
Prompt:
Create 10 multiple-choice questions on the topic: 
"The Water Cycle" for Grade 7 students.

For each question:
- Write one clear question
- Provide 4 options (A, B, C, D)
- Make sure the 3 wrong options are plausible 
  (common misconceptions, not obviously silly)
- Mark the correct answer
- Add a one-line explanation of why it is correct

Difficulty: mix of easy (4), medium (4), hard (2).
\end{lstlisting}



\end{frame}


%%%%%%%%%%%%%%%%%%%%%%%%%%%%%%%%%%%%%%%%%%%%%%%%%%%%%%%%%%%
\begin{frame}[fragile]\frametitle{Explaining Difficult Concepts at Different Levels}

    \begin{itemize}
        \item \textbf{Situation:} You need to explain the same concept to different audiences, a student, a parent, and the school principal.
        \item \textbf{Why this is powerful:} You write the concept once; the AI adapts the vocabulary, analogies, and tone for each audience automatically.
        \item \textbf{Office use:} Same technique for explaining a technical issue to your manager vs. to an IT colleague vs. to a client.
    \end{itemize}

\begin{lstlisting}[language=HTML]
Prompt:
Explain the concept of "compound interest" three times:
1. To a 12-year-old student using a pocket money analogy
2. To a parent at a school financial literacy evening, 
   practical, relatable, no jargon
3. To a school board member in one precise paragraph 
   for an annual report

Keep each explanation self-contained and appropriate 
for the audience.
\end{lstlisting}


\end{frame}

%%%%%%%%%%%%%%%%%%%%%%%%%%%%%%%%%%%%%%%%%%%%%%%%%%%%%%%%%%%
\begin{frame}[fragile]\frametitle{Creating Training and Onboarding Materials}

    \begin{itemize}
        \item \textbf{Situation:} A new employee or student is joining. You need to create a guide, checklist, or FAQ quickly.
        \item \textbf{Time saved:} What normally takes a senior staff member an afternoon of writing now takes 2 minutes to draft, 10 minutes to review and customise.
    \end{itemize}
		
\begin{lstlisting}[language=HTML]
Prompt:
Create a "First Week Onboarding Checklist" for a new 
primary school teacher joining mid-year.

Include:
- Day 1: Essential admin tasks (HR forms, ID card, etc.)
- Day 2-3: Classroom setup and resource access
- Day 4-5: Meeting key people (checklist format)
- One-page FAQ answering the 8 most common questions 
  a new teacher asks (e.g., how to raise a discipline 
  issue, how to request supplies, marking policy)
- A "Who to Call" quick reference card

Format everything so it can be printed on 2 A4 pages.
\end{lstlisting}



\end{frame}


%%%%%%%%%%%%%%%%%%%%%%%%%%%%%%%%%%%%%%%%%%%%%%%%%%%%%%%%%%%
\begin{frame}[fragile]\frametitle{Analysing Survey or Feedback Data}

    \begin{itemize}
        \item \textbf{Situation:} You collected feedback forms, from students, parents, or clients, and need to make sense of 50 free-text responses.
        \item \textbf{What would take hours} of reading, categorising, and writing is done in seconds. You validate and add context, the AI does the heavy lifting.
    \end{itemize}
		
\begin{lstlisting}[language=HTML]
Prompt:
Below are 50 responses from our annual parent 
satisfaction survey. Read all of them carefully.

"""
[Paste all responses here]
"""

Then provide:
1. Top 5 themes that appear most frequently 
   (with example quotes for each)
2. Top 3 specific complaints (most actionable)
3. Top 3 specific compliments (worth celebrating)
4. One surprising or unexpected insight
5. A short paragraph I can include in my report 
   summarising overall sentiment (positive/neutral/
   negative split as a rough percentage)
\end{lstlisting}



\end{frame}

%%%%%%%%%%%%%%%%%%%%%%%%%%%%%%%%%%%%%%%%%%%%%%%%%%%%%%%%%%%
\begin{frame}[fragile]\frametitle{Writing Personalised Student Feedback}

    \begin{itemize}
        \item \textbf{Situation:} You have 30 assignments to mark and need to write meaningful, individual feedback for each student, not a copy-paste.
        \item \textbf{Key technique:} Give the AI the student's actual work and the rubric. It drafts feedback; you personalise and sign off.
		\item \textbf{Time saved:} 2--3 minutes per student instead of 10--15. Always review.% before sharing.

    \end{itemize}

\begin{lstlisting}[language=HTML]
Prompt:
You are a supportive teacher giving written feedback.

Student name: Priya
Assignment: Persuasive essay on school uniforms
Grade level: 8th standard

Rubric criteria: Argument clarity (5), Evidence (5), 
Grammar (5), Structure (5)

Student's essay: """[paste essay here]"""

Write 3-4 sentences of encouraging, constructive 
feedback. Mention one specific strength and one 
clear area for improvement. Avoid generic praise.
\end{lstlisting}



\end{frame}


%%%%%%%%%%%%%%%%%%%%%%%%%%%%%%%%%%%%%%%%%%%%%%%%%%%%%%%%%%%
\begin{frame}[fragile]\frametitle{Building a Personal ``Prompt Library''}

\textbf{The highest-leverage habit you can build:}

\medskip
Once you find a prompt that works well for a recurring task, \textbf{save it as a reusable template}.

\medskip
Replace the \texttt{[PLACEHOLDERS]} each time. \textbf{Your prompt library becomes your personal AI productivity toolkit.}

\begin{lstlisting}[language=HTML]
Example saved prompt templates:

[EMAIL_DELAY]
Write a professional email informing [CLIENT] that 
[PROJECT] is delayed by [DURATION] due to [REASON]. 
Tone: apologetic but confident. Under 150 words.

[LESSON_PLAN]  
Create a 45-minute lesson plan for [TOPIC], Grade [X], 
with prior knowledge of [Y]. Include warm-up, activity, 
5 practice problems, and a closing check question.

[MEETING_SUMMARY]
From these rough notes [NOTES], write: a 3-para summary, 
an action table (Action | Owner | Deadline), and a 
4-line email to attendees.
\end{lstlisting}



\end{frame}


%%%%%%%%%%%%%%%%%%%%%%%%%%%%%%%%%%%%%%%%%%%%%%%%%%%%%%%%%%%
\begin{frame}[fragile]\frametitle{The Golden Rules for Everyday Prompting}

\begin{columns}
    \begin{column}[T]{0.5\linewidth}
    \textbf{Always include:}
    \begin{itemize}
        \item \textbf{Context}, who you are, what the situation is
        \item \textbf{Audience}, who will read the output
        \item \textbf{Tone}, formal / friendly / assertive
        \item \textbf{Format}, bullet list / table / email / paragraph
        \item \textbf{Length}, word count or ``under X sentences''
        \item \textbf{Constraints}, what to avoid
    \end{itemize}
    \end{column}
    \begin{column}[T]{0.5\linewidth}
    \textbf{Always remember:}
    \begin{itemize}
        \item \textbf{First draft, not final output}, always review and personalise
        \item \textbf{Iterate}, if the first output is not quite right, refine the prompt and try again
        \item \textbf{Paste real content}, the AI works best when you give it the actual text, email, or notes to work with
        \item \textbf{Never paste passwords, personal IDs, or confidential data} into public AI tools
        \item \textbf{Save what works}, build your prompt library
    \end{itemize}
    \end{column}
\end{columns}

\medskip
\begin{center}
\textbf{The best prompt is specific, contextual, and tells the AI exactly what ``done'' looks like.}
\end{center}

\end{frame}
